\documentclass{syllabus}

\begin{document}

\thispagestyle{fancy}
\date{Spring 2025}
\author{}
\title{PHY 4200: Quantum Mechanics}
\maketitle



\section*{General Information}

\raisebox{\height}{

\begin{tabular}{lll}
\textbf{Instructor:}    & \multicolumn{2}{l}{Joseph Anderson}                             \\
\textbf{}               & \multicolumn{2}{l}{\href{mailto:janderson19@carthage.edu}{janderson19@carthage.edu}}                    \\
\textbf{}               & \multicolumn{2}{l}{DSC 082}                    \\
\textbf{}               & \multicolumn{2}{l}{(262)551-5871}                               \\ \\
\textbf{Office hours:}  & Tues., Thurs.               & 11 AM - 12 PM                    \\
\textbf{}               & Mon, Wed. Fri                & 2:40-3:30 PM               \\ \\
\textbf{Textbook:}      & \multicolumn{2}{l}{\textit{An Introduction to Quantum Mechanics 3/e}} \\
\textbf{}               & \multicolumn{2}{l}{Griffiths and Schroeter (Cambridge, 2018)}     \\ \\
\textbf{Prerequisites:} & \multicolumn{2}{l}{PHY 2210 (gen. phys. II)}                    \\
\textbf{}               & \multicolumn{2}{l}{MTH 2020 (dif. eq.)}            
\end{tabular}
%
}
\hspace{\fill}
\includegraphics[height=2.5in]{Assets/quantumCover.png}
\hspace{.25in}

\section*{Course overview}

A study of the principles of quantum mechanics. Schrodinger theory and operator algebra are applied to the study of such problems as potential wells and barriers, tunneling, the harmonic oscillator, and the hydrogen atom.

Quantum mechanics is a theory of matter on the smallest scales that upended many of our intuitions about how physics works. It owes its birth to the great height that analytical mechanics (e.g. Hamiltonian mechanics) on the one hand and mathematical analysis (the rise of modern formal mathematics as a discipline) achieved on the other. These combined to answer great conundrums that arose in the theory of atomic matter. Because of this, quantum mechanics (somewhat rightly) has a reputation for being abstract and mathematically complex. In this course, you will learn the advanced mathematics (most especially spectral methods in ordinary and partial differential equations) needed to speak accurately about quantum mechanics and apply that mathematics to the analysis of atomic systems. Topics in this course include:

\begin{enumerate}[label=\textcolor{CarthageDark}{\bfseries\arabic*.}]
    \item An introduction to the wave function, the statistical interpretation, and the momentum operator.
\item An exploration of the time-independent Schrödinger equation in one dimension for a variety of instructive potentials (free particle, infinite square well, harmonic oscillator, delta-function, and finite square well).
\item Function spaces and the theory of Hermitian operators.
\item An analysis of important three dimensional problems, including the hydrogen atom and the treatment of angular momentum/spin.
\item An introduction to multi-electron systems (non-hydrogenic atoms, solids).
\item Variational methods and the role of action principles in deriving the Schrödinger equation.
\end{enumerate}

\subsection*{Student learning outcomes} 
Upon successful completion of this course, the student will be able to:
\begin{enumerate}[label=\textcolor{CarthageDark}{\bfseries\arabic*.}]
\item Discuss wave-particle duality and its role in quantum phenomena.
\item Calculate the stationary states of the Schrödinger equation for some standard potentials.
\item Apply the idea of a stationary state to predict the time evolution of a wave function.
\item Use the wave function to calculate expectation values of observable operators.
\item Understand the relationship of atomic quantum numbers to eigenstates of particular operators in three dimensions and their corresponding wave functions.
\item Be comfortable with vector algebra in infinite-dimensional function spaces.
\item Comment insightfully on correspondence principles which allow for an interface between quantum and classical physics.
\end{enumerate}

\section*{Course components}

\subsection*{Readings}
and their schedule are shown in the Course Schedule below. Griffiths (now with Schroeter) is a classic introductory text and is quite readable. Quantum mechanics is a challenging subject and it is essential that you read the assigned text material before class and come prepared to discuss the assigned reading.

\subsection{Homework} 
Homework exercises will be assigned on a daily basis and due once a week. You are encouraged to discuss them with fellow students and work in small groups: However, rote copying of someone else’s homework is prohibited. Late work will be accepted until solutions are posted (~1 week after due date) and will be penalized 30\%. 

\subsection{Quizzes} 
Short quizzes on the content of each homework set will be administered on the day that homework set is due. Quizzes are mostly on a pass/fail system and are designed to test individual understanding of the material. 

\subsection{Exams} 
Two cumulative mid-term exams will be administered during the semester. The format will be similar to homework problems (thought provoking questions inviting you to explore your understanding) but suitably simplified to fit in a shorter time window. A comprehensive final exam will take place during the final exam period. You will be able to use a single handwritten notes sheet during the exams.

\section*{Grading Scheme}
\subsection*{Component Weights}

\noindent\begin{tabular}{ll}
\textbf{Component} & \textbf{Weight}       \\ \hline
\textbf{Homework}            & 35\%       \\
\textbf{Quizzes}            & 15\%       \\
\textbf{Exams} (2) & 30\% \\
\textbf{Final Exam} & 20\% \\
\end{tabular}%



\lettergrades{a) number of extra credit assignments adequately completed, b) a pattern of seeking help in the course at office hours, c) number of missing homework assignments.}

\tipsforsuccess{upper}

\section*{Policies}
\subsection{General course policies}
Expect a response within one business day for emails sent to the instructor. To ensure a prompt response, format the subject line as: [PHY 4200] *insert subject here*

All required material will be covered in lecture, but can also be found in the parallel readings, so in the case of absence students should consult the scheduled readings. Students found sleeping will be brought to the attention of the class for public ridicule and no effort made to wake them up.

Students are expected to arrive and leave on time. I will attempt to arrive five minutes early and be available for at least five minutes after class.

\collegepolicies

\end{document}
